\documentclass[12pt,a4paper]{ltjsarticle}
\usepackage{amsmath,amssymb}
\usepackage{graphicx}
\usepackage{hyperref}
\usepackage{geometry}
\usepackage{listings}
\usepackage{xcolor}
\usepackage{bm}

\geometry{margin=25mm}

\lstset{
  basicstyle=\ttfamily\small,
  keywordstyle=\color{blue},
  commentstyle=\color{green!60!black},
  stringstyle=\color{orange},
  numbers=left,
  numberstyle=\tiny\color{gray},
  frame=single,
  breaklines=true,
  language=Python
}

\title{\Large 課題1：2次元強磁性イジング模型の転移温度\\
\large 古典モンテカルロ法による数値計算}
\author{}
\date{}

\begin{document}
\maketitle

\tableofcontents
\newpage

%=============================================================
\section{問題設定}
%=============================================================

\subsection{モデルのハミルトニアン}

2次元正方格子上の強磁性イジング模型のハミルトニアンは以下で与えられる：

\begin{equation}
  H = -J \sum_{\langle i,j \rangle} s_i s_j
  \label{eq:hamiltonian}
\end{equation}

ここで，$s_i = \pm 1$ は各格子点 $i$ のスピン変数，$J > 0$ は強磁性結合定数，
$\langle i,j \rangle$ は最近接サイトのペアを表す．

\subsection{系のサイズと境界条件}

系は $L \times L$ の正方格子とし，周期境界条件を課す：
\begin{equation}
  s_{L+1,\, j} = s_{1,\, j}, \qquad s_{i,\, L+1} = s_{i,\, 1}.
\end{equation}
これにより格子はトーラス（ドーナツ面）と同相になり，表面効果を取り除くことができる．
表面効果とは，開放端を持つ格子において端のスピンが内部より隣接数が少ないために生じる
不均一性のことであり，有限サイズ効果と混同される原因となる．

\paragraph{Onsager解との境界条件の比較．}
Onsager（1944）の厳密解では，両方向に周期境界条件を課したトーラスではなく，
横方向（$n$スピン）のみをリング状につないだ\textbf{円筒（cylinder）境界条件}が用いられている
（原論文 p.\,117：``a long strip crystal of finite width ($n$ atoms), joined straight to itself around a cylinder''）．
縦方向は行数を無限大にする熱力学極限で処理される．
本課題のモンテカルロ計算では両方向に周期境界条件（トーラス）を課すが，
熱力学極限 $L\to\infty$ では両者は同じ $T_c$ を与えるため，厳密値との比較に問題はない．

\subsection{厳密解（Onsager解）}

2次元イジング模型はOnsager（1944）により厳密に解かれており，転移温度の厳密値は
\begin{equation}
  \frac{k_B T_c}{J} = \frac{2}{\ln(1+\sqrt{2})} \approx 2.2692
  \label{eq:onsager}
\end{equation}
で与えられる．この式は原論文の付録 式(8.1)：$\sinh(2J/k_BT_c)\sinh(2J'/k_BT_c)=1$
において正方格子（$J=J'$）の場合に $\sinh(2J/k_BT_c)=1$ を解くことで得られる．
本課題ではこの厳密値との比較によりモンテカルロ計算の精度を検証する．

%=============================================================
\section{モンテカルロ法}
%=============================================================

\subsection{メトロポリスアルゴリズム}

カノニカル集団（温度 $T$ 固定）でのサンプリングにはメトロポリスアルゴリズムを用いる．
1回の更新試行は以下のステップからなる：

\begin{enumerate}
  \item サイト $(i,j)$ をランダムに選ぶ．
  \item スピン反転 $s_{i,j} \to -s_{i,j}$ に伴うエネルギー変化
        \begin{equation}
          \Delta E = 2 J s_{i,j} \sum_{\delta} s_{i+\delta}
        \end{equation}
        を計算する（$\delta$ は4つの最近接方向）．
  \item 受理確率
        \begin{equation}
          P_{\rm accept} = \min\!\left(1,\; e^{-\Delta E / k_B T}\right)
        \end{equation}
        でスピン反転を受理または棄却する．
        $\Delta E < 0$（エネルギーが下がる）場合は必ず反転し，
        $\Delta E > 0$（エネルギーが上がる）場合は確率 $e^{-\Delta E/k_BT}$ で反転する．
        高温では上り坂も頻繁に受理されてランダムに振る舞い，
        低温ではほぼ棄却されて系が秩序状態に留まる．
\end{enumerate}

この手続き全体の目的は，温度 $T$ のカノニカル分布
$P(\{s\}) \propto e^{-H(\{s\})/k_BT}$（ボルツマン分布）に従って配置空間をサンプリングすることである．
受理確率がこの形をとるのは，\textbf{詳細釣り合い条件}（detailed balance）
\begin{equation}
  P(\{s\})\, P_{\rm accept}(\{s\}\to\{s'\})
  = P(\{s'\})\, P_{\rm accept}(\{s'\}\to\{s\})
\end{equation}
を満たすように設計されているためであり，これによって十分な時間後に系がボルツマン分布へ収束することが保証される．

\subsection{1モンテカルロステップ（1 MCS）}

1モンテカルロステップ（1 MCS）とは，全 $L^2$ サイトの状態更新を1回ずつ試みることに対応する．
1サイトの更新コストは最近接4サイトの参照のみであるから $O(1)$，
1 MCS 全体のコストは
\begin{equation}
  \text{1 MCS のコスト} = O(L^2).
\end{equation}
したがって，$N_{\rm step}$ ステップの全計算時間は $O(N_{\rm step} \cdot L^2)$ でスケールする．

\subsection{熱化と測定}

十分長い熱化（thermalization）ステップの後，物理量を測定する．
熱化とは系が初期配置の影響を忘れ，平衡分布に収束するまでの過程である．
本計算では $N_{\rm therm}$ ステップ熱化した後，$N_{\rm meas}$ ステップにわたり測定を行う．

%=============================================================
\section{観測量}
%=============================================================

\subsection{磁化と磁化率}

単位体積あたりの磁化および磁化率を以下で定義する：

\begin{align}
  m &= \frac{1}{L^2} \sum_{i} s_i, \\
  \chi &= \frac{L^2}{k_B T}\Bigl(\langle m^2 \rangle - \langle |m| \rangle^2\Bigr).
\end{align}

磁化率 $\chi$ は相転移点付近で発散し，有限系では $L$ に依存したピークを示す．

\subsection{ビンダーキュムラント}

有限サイズスケーリングに用いるビンダーキュムラント（Binder cumulant）を
\begin{equation}
  U_L = 1 - \frac{\langle m^4 \rangle}{3\langle m^2 \rangle^2}
  \label{eq:binder}
\end{equation}
と定義する．$U_L$ は純粋な $\mathbb{Z}_2$ 対称性が破れた秩序相では $U_L \to 2/3$，
無秩序相では $U_L \to 0$ に漸近する．

\subsection{有限サイズスケーリングとビンダークロス}

\paragraph{なぜ有限サイズスケーリングが必要か．}
無限大の系では相転移はちょうど $T_c$ で起きるが，計算機で扱える $L \times L$ の有限系では
転移が温度方向に「ぼける」．$L$ が大きいほどぼけは小さくなるため，
複数の $L$ で計算し，$L \to \infty$ への外挿によって $T_c$ を決定する必要がある．

\paragraph{スケーリング仮説．}
転移点付近では相関長 $\xi \sim |T - T_c|^{-\nu}$ が発散する（$\nu$：臨界指数）．
有限系では $L$ と $\xi$ の比が振る舞いを決めるため，
一般の物理量 $\mathcal{O}$ は
\begin{equation}
  \mathcal{O}(T, L) = L^{p/\nu}\, \tilde{f}\!\left[(T - T_c)\,L^{1/\nu}\right]
\end{equation}
というスケーリング形式をとる．磁化率 $\chi$ などは $L$ 依存の係数 $L^{p/\nu}$ がつくため，
異なる $L$ の曲線を直接比較しにくい．

\paragraph{Fisher--Barber論文との対応．}
Fisher and Barber（1972）の原論文が扱う系は厳密には
「1方向のみ有限な薄膜（$n$ 層 $\times$ 無限大）」であり，
本課題のような両方向が有限な $L \times L$ 系とは設定が異なる．
ただし同論文は周期境界条件のもとでこのスケーリング形式が指数関数的に速く収束することを示しており，
$L \times L$ 系への適用を正当化している（同論文 p.\,1518）．
この枠組みを $L \times L$ 系に適用したのはBinderら以降の研究による．

\paragraph{ビンダーキュムラントの利点．}
ビンダーキュムラント $U_L$（式\eqref{eq:binder}）は $\langle m^4\rangle$ と $\langle m^2\rangle^2$ の比であるため
$L$ 依存の係数が打ち消されて $p = 0$ となり，
\begin{equation}
  U_L(T) = \tilde{f}\!\left[(T - T_c) L^{1/\nu}\right]
\end{equation}
となる．したがって\textbf{異なる $L$ の $U_L(T)$ 曲線は $T = T_c$ で交差する（ビンダークロス）}，
この交点から $T_c$ を直接読み取ることができる（図\ref{fig:binder_schematic}参照）．

\paragraph{2次元イジング模型の臨界指数．}
Onsagerの厳密解から以下の臨界指数が得られている：
\begin{equation}
  \nu = 1,\quad \beta = \frac{1}{8},\quad \gamma = \frac{7}{4}.
\end{equation}
$\nu = 1$ であるから，スケーリング変数は $(T - T_c)L$ とシンプルになる．

\begin{figure}[h]
  \centering
  \begin{minipage}{0.7\textwidth}
    \small
    \textit{（ビンダーキュムラントの模式図：異なる $L$ の曲線が $T_c$ で交差する）}
  \end{minipage}
  \caption{ビンダーキュムラント $U_L$ の温度依存性の模式図．}
  \label{fig:binder_schematic}
\end{figure}

%=============================================================
\section{計算手順のまとめ}
%=============================================================

\begin{enumerate}
  \item $L = 8, 16, 32, 64$ 等，複数のシステムサイズと，$T_c \approx 2.269$ 周辺の温度グリッドを設定する．
        各 $(L,\, T)$ の組み合わせに対して以下を実行する：
        \begin{enumerate}
          \item スピン配置を初期化する（高温ではランダム，低温では全スピン整列など）．
          \item $N_{\rm therm}$ MCS のメトロポリス更新を行い，初期配置の影響を取り除く（熱化）．
                熱化フェーズでも各 MCS でスピン反転の試みは行うが，物理量は記録しない．
                系はこの過程でボルツマン分布へ収束していく．
          \item $N_{\rm meas}$ MCS にわたって各ステップの $m$，$m^2$，$|m|$，$m^4$ を記録する．
                「平衡に達した瞬間に1回だけ記録」では1配置の値しか得られず統計的に無意味である．
                平衡状態に達した後も系はボルツマン分布上をランダムウォークし続けるため，
                その軌跡を多数記録して時間平均を取ることで熱平均
                $\langle \mathcal{O} \rangle = \sum_{\{s\}} \mathcal{O}(\{s\})\,e^{-H/k_BT} / Z$
                を近似する．
          \item 熱平均 $\langle m \rangle$，$\langle m^2 \rangle$，$\langle |m| \rangle$，$\langle m^4 \rangle$ を計算する．
          \item これらから磁化率 $\chi$ およびビンダーキュムラント $U_L$ を計算する．
        \end{enumerate}
  \item 各 $L$ について $\langle|m|\rangle$，$\chi$，$U_L$ を温度 $T$ の関数としてプロットする．
  \item ビンダーキュムラントの交差点から $T_c$ を推定し，Onsager解（式\eqref{eq:onsager}）と比較する．
  \item 計算時間の $L$ 依存性を測定し，$O(L^2)$ スケーリングを確認する．
  \item 小さな $L$ に対するユニットテストを実施する（エネルギー計算，詳細釣り合い条件等）．
\end{enumerate}

%=============================================================
\section{実装上の注意点}
%=============================================================

\begin{itemize}
  \item \textbf{周期境界条件}：インデックス計算に \texttt{mod} 演算を用いる．
  \item \textbf{初期配置}：高温ではランダム配置，低温では全スピン整列配置が収束を速める
        （あるいは両方から始めて一致を確認する）．
  \item \textbf{統計誤差}：独立なサンプルを得るために自己相関時間を考慮する．
        簡便には jackknife 法や bootstrap 法で誤差を推定する．
  \item \textbf{乱数}：メルセンヌ・ツイスタ等の品質の高い擬似乱数生成器を用いる．
  \item \textbf{$O(L^2)$ スケーリングの確認}：$L$ を2倍にするごとに計算時間が約4倍になることを確認する．
\end{itemize}

%=============================================================
\appendix
\section{参考文献}
%=============================================================

\begin{itemize}
  \item L. Onsager, \textit{Phys. Rev.} \textbf{65}, 117 (1944).
  \item M. E. Fisher and M. N. Barber, \textit{Phys. Rev. Lett.} \textbf{28}, 1516 (1972)．（有限サイズスケーリング理論の原論文）
  \item K. Binder, \textit{Z. Phys. B} \textbf{43}, 119 (1981)．（ビンダーキュムラントの導入）
  \item N. Metropolis et al., \textit{J. Chem. Phys.} \textbf{21}, 1087 (1953)．（メトロポリス法）
  \item W. K. Hastings, \textit{Biometrika} \textbf{57}, 97 (1970)．（詳細釣り合い条件の数学的定式化；Metropolis--Hastings法）
  \item M. E. J. Newman and G. T. Barkema,
        \textit{Monte Carlo Methods in Statistical Physics}, Oxford (1999).
\end{itemize}

\end{document}
